\maketitle
\thispagestyle{empty}

\begin{abstract}
LM-IDNet is a lightweight, unsupervised IoT anomaly detector that learns normal
device traffic from TCP, ordinary UDP, SSDP, and ARP packet counts using a
Dirichlet--multinomial model. Separate benign calibration supports likelihood
scoring with compact volume and ARP checks. The Languasco--Migliardi numerical
backend, guarded adaptation, and traceable alert records form a reproducible
pipeline without attack signatures or neural models. This report describes
its statistical basis, data preparation, static and adaptive detectors,
input safeguards, saved state, streaming runtime, simulated edge deployment,
tests, development/final evaluations, and bounded future research.
It retains the earlier experiments, including unsuccessful ones, while
identifying which implementation and adaptation policy produced each result.
The reference software now supports conservative updates, restoration and
rollback, and bounded supplied-window operation in a local Docker simulation.
The completed final-partition evaluation covers 1,008 labelled Chromecast and
Samsung camera windows, with 82.35\% precision, 77.78\% recall, and 80.00\%
pooled F1 under static and both conservative adaptive modes. No conservative
update is promoted on those sequences. The detector retains its DM foundation
with compact count-based scoring and fixed references during updates.
The report is the completed technical and experimental record; independently
unseen external traffic and physical gateway validation remain future work.
\end{abstract}

\vfill
\begin{center}
\begin{minipage}{0.92\textwidth}
\small
\textbf{Evidence boundary, 19 October 2026.}
The selected detector/scoring method is recorded at \texttt{99bc189} on
\texttt{feature/volume-anomaly}; the closing protocol records the current
source and input SHA-256 fingerprints. Historical adaptation, numerical, edge,
and temporal measurements retain their recorded implementation boundaries.
The selected ten-minute models and thresholds are frozen before final replay;
full events, update attempts, metrics, and verification are retained.
Fitting and calibration use normal traffic; final labels are used only for
offline metrics, with detector settings fixed throughout the final run.
The duration study retains its own fitted models, decisions, hashes, and saved
figure. The completed software, final evaluation, and report close the agreed
supplied-window simulation scope.
\end{minipage}
\end{center}

\newpage
\tableofcontents
\newpage
